\documentclass{report}
\usepackage{amsmath,amssymb}
\setlength{\parindent}{0mm}
\setlength{\parskip}{1em}
\begin{document}
\begin{center}
\rule{6in}{1pt} \
{\large Aditi Ghai \\
{\bf An Empirical Comparison of Krylov Subspace Methods for Nonsymmetric Linear Systems from PDEs}}

Department of Applied Math \\ Stony Brook University \\ Stony Brook \\ NY 11794
\\
{\tt aditi.ghai@stonybrook.edu}\\
Aditi Ghai\\
Cao Lu\\
Xiangmin Jiao\end{center}

\providecommand{\tabularnewline}{\\}
Krylov subspace methods are widely used in solving partial differential
equations and large sparse linear systems. For symmetric systems,
CG and MINRES are almost
always the optimal choices as efficient solvers. However, for nonsymmetric
systems, such as those arising from convection-diffusion equations,
the best choice of solvers is much less clearer. Various methods have
been proposed for nonsymmetric systems over the years. Notable methods
include GMRES \cite{gmres}, QMR \cite{qmr}, TFQMR \cite{tfqmr},
Bi-CGSTAB \cite{BSG}, QMRCGSTAB \cite{qcsb}, etc. Unlike for symmetric
case, there is no overall best method, and there is a need of guidelines
in choosing these different types methods. The purpose of this study
is to perform a systematic comparison of the various methods for linear
systems arising from PDE discretizations. We focus on two classes
of methods that are based on either nonsymmetric Lanczos iterations
or Arnoldi iterations. We analyze the algorithm complexity and present
numerical performances of these methods with and without preconditioners,
for linear systems from finite difference or finite element discretizations
in $2$-D and $3$-D.

In this study, we first report some theoretical comparisons of the
methods with three-term recurrences, including QMR, TFQMR, Bi-CGSTAB
and QMRCGSTAB, as well as the methods based on $n$-term recurrence,
such as GMRES. Table~\ref{tab:Comparison-of-operations} compares
the operation counts per iteration and the storage requirements for
these methods, which augment the comparison in \cite{BBC94Templates}
with TFQMR and QMRCGSTAB. There is a tradeoff depending on the average
number of nonzeros in the matrix versus the number of iterations.
We present simple performance models for different classes of linear
systems from PDE discretizations based on finite difference or finite
element discretizations in 2-D and 3-D to provide practical guidelines
in choosing these different Krylov subspace solvers.

\begin{table}
\caption{\label{tab:Comparison-of-operations}Comparison of operations per
iteration and memory requirements of various methods. $m$ denotes
the number of rows, $n$ denotes the average number of nonzeros per
row, and $k$ denotes the iteration count.}


\centering{}%
\begin{tabular}{c|c|c|c|c|c}
\hline
& MatVec & & Inner & & \tabularnewline
& Prod. & DAXPY & Prod. & FLOPs & Storage\tabularnewline
\hline
\hline
QMR & 2 & 12 & 2 & $4m(n+7)$ & $m(n+16)$\tabularnewline
\hline
TFQMR & 2 & 10 & 4 & $4m(n+7)$ & $m(n+8)$\tabularnewline
\hline
Bi-CGSTAB & 2 & 6 & 4 & $4m(n+5)$ & $m(n+10)$\tabularnewline
\hline
QMRCGSTAB & 2 & 8 & 6 & $4m(n+7)$ & $m(n+13)$\tabularnewline
\hline
GMRES & 1 & $k$+1 & $k$+1 & $2m(n+2k+2)$ & $m(n+k+5)$\tabularnewline
\hline
\end{tabular}
\end{table}


To compliment the theoretical analysis, we also report some empirical
comparisons of the different methods in terms of their convergence
rates with and without preconditioners. We evaluate three preconditioners,
including incomplete LU, Gauss Seidel, and Chebyshev polynomials.
From the results, we observe that methods GMRES is more robust, with
monotonically decreasing residuals, QMRCGSTAB, QMR, and TFQMR are
less robust with nearly monotonically decreasing residual, and Bi-CGSTAB
is the most efficient but also the least robust, with oscillatory
residuals. For these different methods, we observe different accelerations
with different preconditioners for different classes of problems.
These results help provide some practical guidelines in choosing different
preconditioners for different types of linear systems, and also motivate
the development of hybrid solvers for sparse linear systems arising
from PDE discretizations.

\begin{thebibliography}{1}

\bibitem{BBC94Templates}
R.~Barrett, M.~Berry, T.~F. Chan, J.~Demmel, J.~Donato, J.~Dongarra,
V.~Eijkhout, R.~Pozo, C.~Romine, and H.~{Van der Vorst}.
\newblock {\em Templates for the Solution of Linear Systems: Building Blocks
for Iterative Methods}.
\newblock SIAM, Philadelphia, PA, 2nd edition, 1994.

\bibitem{qcsb}
T.~F. Chan, E.~Gallopoulos, V.~Simoncini, T.~Szeto, and C.~H. Tong.
\newblock A quasi-minimal residual variant of the {Bi-CGSTAB} algorithm for non
symmetric systems.
\newblock {\em SIAM J. Sci. Comput.}, 15(2):338--347, 1994.

\bibitem{tfqmr}
R.~W. Freund.
\newblock A transpose-free quasi-minimal residual algorithm for non-hermitian
linear systems.
\newblock {\em SIAM J. Sci. Comput.}, 14(2):470--482, 1993.

\bibitem{qmr}
R.~W. Freund and N.~M. Nachtigal.
\newblock An implementation of the {QMR} method based on coupled two-term
recurrences.
\newblock {\em SIAM J. Sci. Comput.}, 15(2):313--337, 1994.

\bibitem{gmres}
Y.~Saad and M.~H. Schultz.
\newblock {GMRES}: A generalized minimal residual algorithm for solving non
symmetric linear systems.
\newblock {\em SIAM J. Sci. Comput.}, 7(3):856--869, 1986.

\bibitem{BSG}
H.~A. {Van der Vorst}.
\newblock {BI-CGSTAB}: A fast and smoothly convergent variant of {Bi-CG} for
solution of non symmetric linear systems.
\newblock {\em SIAM J. Sci. Comput.}, 13(2):631--644, 1992.

\end{thebibliography}


\end{document}
